\section{An exact fixed-axis singularity of the arithmetic trace}
\label{sec:cq-ns-axis-arithmetic-jet}

The complete companion fixed-axis proof \cite{cqNSAxis} is included with a local proof of its endpoint summation formula. The angular-momentum transform can see a fixed-axis derivative of the
actual source field directly, without first applying its inverse. The
reason is a specific feature of the construction: every positive-order
azimuthal coefficient has zero axis value. We retain the source amplitude
\(C\), the parameter \(h\), the physical viscosity \(\nu\), all cutoffs and
the original time \(t=1-\tau\). The calculation uses the constructed
field of Section~\ref{sec:cq-NS-actual-witness}; it does not certify the
entire existence argument of the released manuscript independently.

\subsection{The axis trace of the fully corrected velocity}

Here and only here let \(q_{\rm src}(z,t)\) denote the scalar similarity
coordinate called \(q\) in the source, not the material label in
Section~\ref{sec:cq-NS-profile-transfer}. Its source coordinates are
\[
 \tau=1-t=q_{\rm src}(1-\eta^2),\qquad
 z=q_{\rm src}^{D}\eta,\qquad
 X=\frac{r^2}{2q_{\rm src}},\qquad
 D=\frac12-h,\quad A=\frac12+h.
\]
In particular at \(z=0\) one has exactly \(\eta=0\) and
\(q_{\rm src}=\tau\).
Source (5.1), pp.~45--46, writes each azimuthal coefficient as
\[
 u_{\theta,n}
   =\frac rC q_{\rm src}^{-1-h+2nh}\varphi_n(X,\eta).
\]
Its Lemma~5.1, p.~47, prescribes
\(\varphi_n(0,\eta)=0\) for every \(n\ge1\).
The radial extensions and moment corrections are supported away from
the unchanged inner interval. Source Proposition~5.5 and its proof,
pp.~60--62, realize this series with cutoffs \(\chi(c_nq_{\rm src})\);
these depend on \(z,t\), not \(r\), and the realized angular component is
\begin{equation}\label{eq:cq-naj:base-series}
 \frac{u_{\theta,B}}r
 =\frac{q_{\rm src}^{-1-h}}C
 \left(\varphi_0(X,\eta)+
   \sum_{n\ge1}\chi(c_nq_{\rm src})
                       q_{\rm src}^{2nh}\varphi_n(X,\eta)\right).
\end{equation}
For every fixed preterminal \(z,t\), this sum is finite in a neighborhood
of that point: \(c_n\to\infty\), while \(q_{\rm src}>0\) and \(\chi\)
has compact support. Thus evaluation at \(X=0\), and the smooth extension
of the quotient \(u_{\theta,B}/r\), require no exchange of a terminal
limit with an infinite sum.

The leading axis datum is not an unknown positive remainder. Source
(B.3), p.~145, and Proposition~B.2, pp.~146--147, give
\[
 \varphi_0(0,\eta)=\varphi_*(\eta),\qquad
 \varphi_*(\eta)=
  \exp\left(\Lambda\int_0^\eta\zeta_*(w)\,dw\right),\qquad
 \varphi_*(0)=1.
\]
Here \(\zeta_*\) is the source's auxiliary axis function, not the
Riemann zeta function. Neither its value nor \(C\) is rescaled.

All annular wave and mean potentials, including the finite
initialization and the later cutoff-summed stages, vanish near the axis
for each \(t<1\): source Proposition~9.9, (9.21), pp.~114--117.
Only finitely many stages are present locally at such a time. Their
Cartesian derivatives therefore vanish on a common neighborhood of the
axis. The base potential is azimuthal and independent of angle, so its
curl has only radial and axial components. Consequently it contributes
no additional azimuthal term to \eqref{eq:cq-naj:base-series}.
Finally the spatial and temporal cutoffs in source (10.4), p.~118,
are identically one near the origin for all sufficiently small
positive \(\tau\). Their derivatives vanish there as well.
These are facts about the actual summed source field, not a
replacement of that field by its leading approximation
\cite[Lemma 5.1; Proposition 5.5; (9.21); (10.4); (B.3)]{cqNS2026}.

It follows that there is a source-determined \(\tau_1>0\) such that
\begin{equation}\label{eq:cq-naj:axis-star}
 \left.\frac{u_{\theta,*}(r,0,1-\tau)}r\right|_{r=0}
       =\frac1C\tau^{-1-h}\qquad(0<\tau<\tau_1).
\end{equation}
The quotient denotes its proved smooth extension, not division by
zero. For the physical viscosity map
\(u_\nu(x,t)=\sqrt\nu\,u_*(x/\sqrt\nu,t)\), put
\(r_*=r/\sqrt\nu\). Then, for \(r>0\),
\[
 \frac{u_{\theta,\nu}(r,0,t)}r
 =\frac{\sqrt\nu\,u_{\theta,*}(r_*,0,t)}{\sqrt\nu\,r_*}
 =\frac{u_{\theta,*}(r_*,0,t)}{r_*}.
\]
Passing to the smooth axis extension proves the same identity
\eqref{eq:cq-naj:axis-star} for \(u_\nu\). This cancellation follows from the
actual coordinate map; viscosity has not been set to one.

Use the original physical radial coordinate
\(\sigma=r^2/2\) and angular average \( (2\pi)^{-1}\int_0^{2\pi}\).
Let
\[
 B_\nu(\sigma,z,t)=
    \frac1{2\pi}\int_0^{2\pi} r u_{\theta,\nu}(r,\theta,z,t)\,d\theta .
\]
In the axis neighborhood just established the full field is
axisymmetric and \(B_\nu=2\sigma(u_{\theta,\nu}/r)\).
The exact consequence is
\begin{equation}\label{eq:cq-naj:B-axis}
 B_\nu(0,0,1-\tau)=0,\qquad
 \partial_\sigma B_\nu(0,0,1-\tau)=\frac2C\tau^{-1-h}.
\end{equation}
No \(O(\tau^{2h})\) error is present in this first axis derivative:
each positive-order axis value is exactly zero. This does not remove
the actual remainder at the distinct nonzero \(X_{\rm in}\) sample
of Section~\ref{sec:cq-NS-actual-witness}.

\subsection{The endpoint summation identity with its complete remainder}
\label{sec:cq-angular-EM-proof}

The endpoint calculation below uses an exact summation formula.
We give its proof here so that the fixed-axis result has no missing
local dependency. Define the Bernoulli polynomials by the convergent
power series, for the auxiliary variable $v$ near zero,
\[
 \frac{ve^{xv}}{e^v-1}=\sum_{n=0}^\infty B_n(x)\frac{v^n}{n!},
 \qquad B_n=B_n(0).
\]
The singularity at $v=0$ is removable. Differentiation in $x$ gives
$B_n'=nB_{n-1}$; subtracting the series at $x=0$ from that at
$x=1$ gives $B_n(1)=B_n(0)$ for $n\ne1$ and
$B_1(1)-B_1(0)=1$. Expansion to second degree gives
$B_1(x)=x-1/2$ and $B_2=1/6$.
The function $v/(e^v-1)+v/2$ is even, by direct substitution
of $-v$. Therefore $B_{2k+1}=0$ for every $k\geq1$.
Write $\widetilde B_n(y)=B_n(y-\lfloor y\rfloor)$ away
from integers, with the common endpoint value for $n\geq2$.
For these $n$ the function is continuous, bounded and periodic.

Let $f$ be smooth and compactly supported on $[0,\infty)$,
let $x>0$, and let $M\geq1$ be an integer. The exact identity is
\begin{equation}
\begin{split}
 \sum_{n\geq1}f(nx)
 &=\frac1x\int_0^\infty f(v)\,dv-\frac{f(0)}2
 -\sum_{k=1}^M\frac{B_{2k}}{(2k)!}
                         x^{2k-1}f^{(2k-1)}(0)+\mathcal R_M(f,x),\\
 \mathcal R_M(f,x)
 &=-\frac{x^{2M-1}}{(2M)!}
    \int_0^\infty\widetilde B_{2M}(v/x)f^{(2M)}(v)\,dv,\\
 |\mathcal R_M(f,x)|
 &\leq\frac{\|B_{2M}\|_{L^\infty([0,1])}}{(2M)!}
                          x^{2M-1}\|f^{(2M)}\|_{L^1(0,\infty)}.
\end{split}
\label{eq:cq-angular-EM}
\end{equation}
To prove it first take mesh one on $[0,N]$, with integer $N$
beyond the support. On every unit interval integration by parts gives
\[
 \int_n^{n+1} B_1(v-n)f'(v)\,dv
 =\frac{f(n)+f(n+1)}2-\int_n^{n+1}f(v)\,dv.
\]
Summing produces the trapezoidal sum minus the integral.
Repeated integration by parts using $B_j'=jB_{j-1}$ yields
the endpoint terms $B_{2k}(f^{(2k-1)}(N)-f^{(2k-1)}(0))/(2k)!$
and the remaining integral
$-\int_0^N\widetilde B_{2M}f^{(2M)}/(2M)!$.
The internal boundary terms cancel because the endpoint values
of $B_j$ agree for $j\geq2$; the odd endpoint coefficients
above degree one vanish by the already proved parity.
All upper endpoint derivatives vanish by the choice of $N$.
Apply this identity to $v\mapsto f(xv)$ and substitute $u=xv$
in each integral. This gives every power and sign in
\eqref{eq:cq-angular-EM}, including the stated norm bound.

For completeness the value of zeta needed at the first axis pole
also follows from this formula. Apply its finite-interval proof
on $[1,N]$ to $v^{-w}$ and initially let $N\to\infty$
for $\operatorname{Re}w>1$. For $M=2$ the result is
\[
 \zeta(w)=\frac1{w-1}+\frac12
    +\frac{B_2}{2!}w+\frac{B_4}{4!}w(w+1)(w+2)
    -\frac{(w)_4}{4!}\int_1^\infty
                     \widetilde B_4(v)v^{-w-4}\,dv.
\]
Here $(w)_4=w(w+1)(w+2)(w+3)$.
Boundedness of the periodic polynomial proves that the last
integral is holomorphic for $\operatorname{Re}w>-3$, with
all complex derivatives justified by the integrable logarithmic
weights on compact subsets. Thus the equality continues there
apart from the pole at one. Substitution $w=-1$ cancels the
first two terms, and the two higher products vanish. Hence
$\zeta(-1)=-B_2/2=-1/12$, with its sign fixed by the actual
summation formula. These classical formulas agree with
\cite[\S2.10(i)]{cqDLMFEM} and \cite[\S25.6(i)]{cqDLMFZeta}.


\subsection{The arithmetic operator on every endpoint derivative}

We prove the endpoint map on its full smooth compactly supported
domain. Let \(a\in C^\infty([0,\infty))\) have compact support and
\(a(0)=0\). For the same fixed \(h\) define
\[
 \phi_a(x)=a(x)-2^h a(2x),\qquad
 \psi_a(x)=\phi_a(x)-\tfrac12\phi_a(x/2),\qquad
 (\mathcal A_h a)(x)=\sum_{n\ge1}\psi_a(nx)\quad(x>0).
\]
This is exactly the arithmetic operator used on the nonlinear fluxes.
There is no radius change hidden in its argument \(x\).
Every sum and its derivatives are locally finite for \(x>0\).
Substitution in the two integrals proves
\(\int_0^\infty\psi_a=0\), and \(\psi_a(0)=0\).
The full endpoint derivatives of the input are
\begin{equation}\label{eq:cq-naj:psi-jet}
 \psi_a^{(j)}(0)
  =(1-2^{-j-1})(1-2^{h+j})a^{(j)}(0)\qquad(j\ge0).
\end{equation}

For clarity the passage from the sum to actual endpoint derivatives
is justified here, rather than differentiating a formal asymptotic
series. We use the Euler--Maclaurin identity proved with its remainder
in \eqref{eq:cq-angular-EM}; its classical provenance is
\cite[\S2.10(i), (2.10.1)]{cqDLMFEM}. With
\(f_j(v)=v^j\psi_a^{(j)}(v)\), direct differentiation on \(x>0\)
gives
\[
 x^j(\mathcal A_ha)^{(j)}(x)=\sum_{n\ge1}f_j(nx).
\]
All these \(f_j\) are smooth and compactly supported. Integration by
parts \(j\) times gives
\(\int f_j=(-1)^j j!\int\psi_a=0\); the factors \(v^j\) make
each zero-end boundary term vanish. Their axis derivatives are
\[
 f_j^{(r)}(0)=0\quad(r<j),\qquad
 f_j^{(r)}(0)=\frac{r!}{(r-j)!}\psi_a^{(r)}(0)\quad(r\ge j).
\]
Choose the Euler--Maclaurin order \(M\) so that \(2M-1>j\).
After division by \(x^j\), its remainder is bounded by
\[
 \frac{\|B_{2M}\|_\infty}{(2M)!}\,
 x^{2M-1-j}\|f_j^{(2M)}\|_1,
\]
and hence tends to zero. Terms below degree \(j\) vanish, so every
differentiated series has a finite limit at zero. The only constant
term for \(j\ge1\) is the Bernoulli term with \(2k-1=j\); for even
\(j\ge2\) it is absent. The fundamental theorem of calculus on
\([\epsilon,x]\), followed by \(\epsilon\downarrow0\), now identifies
these derivative limits with successive derivatives of the extended
function. Boundedness of the next derivative justifies the passage
under each integral. Thus \(\mathcal A_ha\) is smooth at zero, and
\begin{equation}\label{eq:cq-naj:all-jets}
 \begin{split}
 (\mathcal A_ha)(0)&=0,\\
 (\mathcal A_ha)^{(j)}(0)&=
 -\frac{B_{j+1}}{j+1}
 (1-2^{-j-1})(1-2^{h+j})a^{(j)}(0),\qquad j\ge1.
 \end{split}
\end{equation}
The Bernoulli convention is \(B_2=1/6\) and odd \(B_{j+1}=0\)
for even \(j\ge2\).
In particular
\begin{equation}\label{eq:cq-naj:first-map}
 (\mathcal A_ha)'(0)=
       \frac{2^{h+1}-1}{16}\,a'(0).
\end{equation}
The coefficient is strictly positive for the source \(h>0\).

The preceding proofs apply also after any fixed \(z,t\) derivative of
a smooth compactly supported parameter family with common support on
compact preterminal parameter sets. Every bound then holds uniformly
on those sets, with constants depending on the derivative order.
This proves compatibility of the endpoint map with those parameter
derivatives. It supplies no uniform estimate as \(t\uparrow1\).

\subsection{A directly visible singularity at a fixed arithmetic coordinate}

Apply this same operator to the actual compactly supported
\(a(\sigma)=B_\nu(\sigma,z,t)\). Write its arithmetic image as
\[
 \mathcal B_\nu(x,z,t)=\mathcal A_hB_\nu(x,z,t).
\]
The domain and all nonlinear fluxes are those already given in the
preceding angular-momentum section. Equations \eqref{eq:cq-naj:B-axis} and
\eqref{eq:cq-naj:first-map} now prove the exact formula
\begin{equation}\label{eq:cq-naj:exact-blowup}
 \boxed{\displaystyle
 \partial_x\mathcal B_\nu(0,0,1-\tau)
    =\frac{2^{h+1}-1}{8C}\,\tau^{-1-h},
       \qquad 0<\tau<\tau_1,\quad \nu>0 .}
\end{equation}
Here \(x=0,z=0\) are fixed coordinates, not a shrinking sampling path,
and \(\mathcal B_\nu(0,z,t)=0\) for every preterminal \(t\).
Thus the value is zero while its first \(x\)-derivative diverges
positively. For any fixed \(\epsilon>0\), the \(C^1([0,\epsilon])\)
seminorm is at least the displayed endpoint derivative. In particular
no extension continuous at \(t=1\) with values in
\(C^1([0,\epsilon])\) can agree with this family.
This assertion follows from an exact derivative of the original
arithmetic sum, not merely a lower bound obtained after its inverse.

For every integer \(m\ge0\), with \((1+h)_m=\prod_{j=0}^{m-1}(1+h+j)\) and empty product one, differentiation of the exact equality gives
\[
 \partial_t^m\partial_x\mathcal B_\nu(0,0,t)
  =\frac{2^{h+1}-1}{8C}(1+h)_m(1-t)^{-1-h-m},
 \qquad 1-\tau_1<t<1 .
\]
Indeed \(d_t(1-t)^{-\alpha}=\alpha(1-t)^{-\alpha-1}\).
These are equalities on an open preterminal interval. They make no
assignment of finite terminal derivatives.

The projection onto \(B_\nu\) has not discarded the other velocity
data in the full correspondence. Its exact smooth section and kernel
are
\[
 R B_\nu=\frac{B_\nu(r^2/2,z,t)}{r^2}(-x_2,x_1,0),\qquad
 w_\nu=u_\nu-RB_\nu,\qquad \langle r(w_\nu)_\theta\rangle=0,
\]
where \(B_\nu/\sigma\) has its smooth integral extension at zero.
The preceding section proves the inverse
\((B_\nu,w_\nu)\mapsto RB_\nu+w_\nu\) and retains every covariance
in both momentum fluxes. The signal \eqref{eq:cq-naj:exact-blowup} concerns
the \(B_\nu\) component of that complete map. It does not replace the
full Navier--Stokes equation by a closed scalar equation for \(B_\nu\).

\subsection{The exact Mellin pole carrying this axis signal}

For fixed preterminal parameters, put
\[
 \widehat B_\nu(s)=\int_0^\infty
                  B_\nu(\sigma,z,t)\sigma^{s-1}\,d\sigma,\qquad
 \mathfrak B_\nu(s)=\int_0^\infty
                  \mathcal B_\nu(x,z,t)x^{s-1}\,dx .
\]
Both integrals are initially holomorphic on \(\Re s>-1\).
The full transform calculation in the preceding section gives
\[
 \mathfrak B_\nu(s)=W_h(s)\widehat B_\nu(s),\qquad
 W_h(s)=\zeta(s)(1-2^{s-1})(1-2^{h-s}).
\]
The pole of \(\zeta\) at \(1\) is canceled by the displayed first
dilation factor, making \(W_h\) entire.

Here is the continuation needed at \(s=-1\). For any fixed \(b>0\),
split the integral at \(b\) and subtract the linear Taylor term on
\([0,b]\). Since \(B_\nu(0)=0\), the exact formula is
\[
 \widehat B_\nu(s)=
 \int_0^b [B_\nu(\sigma)-B_{\nu,\sigma}(0)\sigma]\sigma^{s-1}d\sigma
 +\int_b^\infty B_\nu(\sigma)\sigma^{s-1}d\sigma
 +\frac{B_{\nu,\sigma}(0)b^{s+1}}{s+1}.
\]
Taylor's integral remainder is \(O(\sigma^2)\); the first integral
is holomorphic for \(\Re s>-2\), including after any complex
\(s\) derivative on a compact subset, because the additional powers
of \(\log\sigma\) remain integrable. The second is entire by compact
support away from zero. Hence
\(\operatorname{Res}_{s=-1}\widehat B_\nu=B_{\nu,\sigma}(0)\).
The identical argument applies to the proved smooth
\(\mathcal B_\nu\), so its residue is \(\partial_x\mathcal B_\nu(0)\).

The classical value \(\zeta(-1)=-1/12\)
\cite[\S25.6(i)]{cqDLMFZeta} gives
\[
 W_h(-1)=-\frac1{12}\left(1-\frac14\right)(1-2^{h+1})
         =\frac{2^{h+1}-1}{16}>0.
\]
Equality on \(\Re s>-1\) therefore continues to the meromorphic
functions just constructed and agrees with the direct endpoint proof.
At \(z=0,t=1-\tau\) the pole is genuine and its residue is exactly
\begin{equation}\label{eq:cq-naj:residue}
 \operatorname{Res}_{s=-1}\mathfrak B_\nu(s,0,1-\tau)
     =\frac{2^{h+1}-1}{8C}\tau^{-1-h}.
\end{equation}
There is no zero-factor compensation at this pole, since \(W_h(-1)>0\).
The sign of \(\zeta(-1)\) alone is not the sign of the arithmetic
residue: both dilation factors have been evaluated explicitly.

This identifies what the construction propagates in this channel:
an exact axis derivative, equivalently a growing residue at the fixed
Mellin point \(-1\). The identity retains the full
\(\widehat B_\nu\), its flux equations, and its inverse maps.
The nontrivial zeros of \(\zeta\) remain the zero factors specified
in the full product formula; no off-critical zero is inferred from
this axis pole or from its positive residue.
